% 자동 생성 — scripts/build_refs_tex.py. 손으로 고치지 않는다.
% 출처: paper/refs_kr.py (학과 안내서 형식)
\begin{thebibliography}{20}
\setlength{\itemsep}{6pt}
\bibitem{jeon2025}
  전상현, 김성용, 박종선. (2025). “토지거래허가구역 지정이 아파트 가격에 미치는 영향에 관한 연구”, 융합사회와 공공정책, 제19권, 제1호, pp. 169~189.
\bibitem{lee2025}
  이한웅, 이춘원. (2025). “토지거래허가구역 지정이 매매가와 전세가에 미치는 영향: 잠실동 아파트 사례분석”, 대한부동산학회지, 제43권, 제2호, pp. 161~180.
\bibitem{yang2025}
  양지영, 정준호. (2025). “토지거래허가구역 지정의 풍선효과 분석: 잠실·삼성·대치·청담동 사례를 중심으로”, 한국콘텐츠학회논문지, 제25권, 제5호, p. 447.
\bibitem{rambachan2023honest}
  Rambachan, A., \& Roth, J. (2023). “A More Credible Approach to Parallel Trends”, Review of Economic Studies, Vol. 90, No. 5, pp. 2555~2591.
\bibitem{dechaisemartin2020twoway}
  de Chaisemartin, C., \& D’Haultfœuille, X. (2020). “Two-Way Fixed Effects Estimators with Heterogeneous Treatment Effects”, American Economic Review, Vol. 110, No. 9, pp. 2964~2996.
\bibitem{goodmanbacon2021}
  Goodman-Bacon, A. (2021). “Difference-in-Differences with Variation in Treatment Timing”, Journal of Econometrics, Vol. 225, No. 2, pp. 254~277.
\bibitem{callaway2021}
  Callaway, B., \& Sant’Anna, P. H. C. (2021). “Difference-in-Differences with Multiple Time Periods”, Journal of Econometrics, Vol. 225, No. 2, pp. 200~230.
\bibitem{sunabraham2021}
  Sun, L., \& Abraham, S. (2021). “Estimating Dynamic Treatment Effects in Event Studies with Heterogeneous Treatment Effects”, Journal of Econometrics, Vol. 225, No. 2, pp. 175~199.
\bibitem{arkhangelsky2021}
  Arkhangelsky, D., Athey, S., Hirshberg, D. A., Imbens, G. W., \& Wager, S. (2021). “Synthetic Difference-in-Differences”, American Economic Review, Vol. 111, No. 12, pp. 4088~4118.
\bibitem{abadie2010synthetic}
  Abadie, A., Diamond, A., \& Hainmueller, J. (2010). “Synthetic Control Methods for Comparative Case Studies: Estimating the Effect of California’s Tobacco Control Program”, Journal of the American Statistical Association, Vol. 105, No. 490, pp. 493~505.
\bibitem{mackinnon2017wild}
  MacKinnon, J. G., \& Webb, M. D. (2017). “Wild Bootstrap Inference for Wildly Different Cluster Sizes”, Journal of Applied Econometrics, Vol. 32, No. 2, pp. 233~254.
\bibitem{webb2014reworking}
  Webb, M. D. (2014). “Reworking Wild Bootstrap Based Inference for Clustered Errors”, Working Paper No. 1315, Queen’s University Department of Economics, Kingston, Canada.
\bibitem{fisher1935design}
  Fisher, R. A. (1935). “The Design of Experiments”, Oliver and Boyd, Edinburgh.
\bibitem{simonsohn2020specification}
  Simonsohn, U., Simmons, J. P., \& Nelson, L. D. (2020). “Specification Curve Analysis”, Nature Human Behaviour, Vol. 4, No. 11, pp. 1208~1214.
\bibitem{bailey1963regression}
  Bailey, M. J., Muth, R. F., \& Nourse, H. O. (1963). “A Regression Method for Real Estate Price Index Construction”, Journal of the American Statistical Association, Vol. 58, No. 304, pp. 933~942.
\bibitem{caseshiller1987}
  Case, K. E., \& Shiller, R. J. (1987). “Prices of Single-Family Homes Since 1970: New Indexes for Four Cities”, New England Economic Review, pp. 45~56.
\bibitem{bestkleven2018}
  Best, M. C., \& Kleven, H. J. (2018). “Housing Market Responses to Transaction Taxes: Evidence from Notches and Stimulus in the U.K.”, Review of Economic Studies, Vol. 85, No. 1, pp. 157~193.
\bibitem{repo}
  이찬희. (2026). “서울 토지거래허가구역 인과효과 추정 — 코드·자료·결정기록”, GitHub 저장소. \url{https://github.com/cosmicpotato2047/capstone-ltz-causal}
\bibitem{site}
  이찬희. (2026). “연구 결과 공개 페이지 — 주차 보고서·발표 자료·결과 대시보드”, GitHub Pages. \url{https://cosmicpotato2047.github.io/capstone-ltz-causal/}
\bibitem{molit_api}
  국토교통부. “아파트 매매 실거래가 상세 자료(RTMSDataSvcAptTradeDev)”, 공공데이터포털 오픈 API. \url{https://www.data.go.kr/}
\end{thebibliography}
